\section{Picard-rank-one Fano threefolds}
\label{sec:fano-classification}
The cubic is the rank-two model for a larger numerical obstruction.
Five families are detected by the canonical rank-two residue. Four more
have a repeated factor of even rank three; there we count its odd
cohomology. No Hodge group is needed for either construction.

\paragraph{A resonant example: degree two.}
For an index-two Fano threefold of degree two, the small matrix in
Appendix~\ref{sec:fano-data} has $a=48$, $b=160$ in
\eqref{eq:universal-residue}. Substitution gives
\[
 R_{48,160}=\begin{pmatrix}-9/8&1\\-9/64&1/8\end{pmatrix},
 \qquad \det(TI-R_{48,160})=T(T+1).
\]
The residue eigenvalues are $0,-1$. Their exponent classes modulo $\Z$
coincide, so this block contributes zero to $I_{\mathrm{exp}}$.
The canonical lattice retains their nonzero difference:
$\delta^\sharp=1$, so it contributes one to $I_{\mathrm{lat}}$.
The rank-one complements contribute nothing. This is why retaining the
lattice detects a family that the exponent count misses.

\subsection{Odd dimension on rank-three factors}
Use the full quantum cohomology supermodule over the same even bulk base.
The even Frobenius algebra acts on it. Its spectral projectors split both
parities, and the full even rank determines the primary factor.
Define
\[
 O_3(T)=\sum_{\rk E^{\mathrm{ev}}=3}\dim E^{\mathrm{odd}},
\]
where the sum runs over whole primary factors. We use the full odd dimension,
so no divisibility convention is needed.

\begin{lemma}[Odd dimension and surface vanishing]
\label{lem:odd-selector}\coverage{absent}%
An even-rank-one primary factor of a Frobenius superalgebra has zero odd
part. The count $O_3$ satisfies
\[
 O_3(\Bl_ZT)=O_3(T)+(c-1)O_3(Z)
\]
for a smooth codimension-$c$ center. It vanishes in dimensions at most two
and is a birational invariant of smooth projective fourfolds.
\end{lemma}
\begin{proof}
Let $e$ be the unit of an even-rank-one factor. For odd $a,b$ write
$ab=\gamma e$. Supercommutativity gives $(ab)^2=0$, hence $\gamma=0$.
The odd Frobenius pairing is therefore zero; nondegeneracy forces the odd
part to be zero.

The blowup comparison and its inverse preserve parity on the full fiber.
They preserve full even ranks; independent occurrence-unit shifts separate
their primary factors as in Proposition~\ref{prop:numerical-blowup}.
The odd dimensions consequently add, and Tate suspensions preserve parity.
Points and curves have even rank at most two. The ruled-product computation
in Lemma~\ref{lem:ruled-product} has no even-rank-three factor.
Point blowups and surface factorization extend this to all ruled surfaces.
Projective space has simple generic factors.

For a minimal surface with nef canonical class, the degree argument in
Proposition~\ref{prop:atomic-lowdim} gives one whole even factor of rank
$b_2+2$. It is selected only if $b_2=1$. In that case a nonzero holomorphic
one-form $\alpha$ would give a nonzero real class
$[i\alpha\wedge\bar\alpha]$ of square zero: its product with a K\"ahler
class has positive integral. This is impossible when $H^2$ is generated
by an ample class. Thus $H^1=0$ and, by duality, $H^3=0$.
The selected odd dimension is zero. Surface classification and point
blowups finish the vanishing proof. Weak factorization in dimension four
then has zero correction at every step.
\uses{prop:numerical-blowup,prop:atomic-lowdim,lem:ruled-product}%
\imports{IritaniBlowup:thm-5.18,AKMW:weak-factorization,BeauvilleSurfaces:minimal-surface-classification}%
\end{proof}

\subsection{The nine detected families}
For index two put $d=H^3$, where $-K_X=2H$. For index one put
$g=\tfrac12(-K_X)^3+1$. Let $\mathcal F$ be the nine families in the
following table. The even cohomology of each member has rank four, and
its only odd cohomology is $H^3$.

\Needspace{20\baselineskip}
\begin{proposition}[Fano endpoint data]
\label{prop:fano-endpoints}\coverage{absent}%
Every smooth member of the indicated family has the following repeated
small even primary factor; all complementary factors have rank one.
The repeated factor is cyclic. The last two columns give the generic
counts after continuation.
\begin{center}
\begin{tabular}{@{}lcccc@{}}
\toprule
Family & Repeated even rank & $\delta^\sharp$ & $I_{\rm lat}$ & $O_3$\\
\midrule
Index two, $d=1$ &2&$16/9$&1&0\\
Index two, $d=2$ &2&$1$&1&0\\
Index two, $d=3$ &2&$4/9$&1&0\\
Index one, $g=2$ &3&---&0&104\\
Index one, $g=3$ &3&---&0&60\\
Index one, $g=4$ &3&---&0&40\\
Index one, $g=5$ &3&---&0&28\\
Index one, $g=6$ &2&$1$&1&0\\
Index one, $g=8$ &2&$4/9$&1&0\\
\bottomrule
\end{tabular}
\end{center}
\end{proposition}
\begin{proof}
Appendix~\ref{sec:fano-data} gives the nine matrices, their source
normalizations and the exact finite calculation. The degree-one, two and
three residue computation is also Proposition~\ref{prop:universal-rank-two-residue}.
The family formulas apply to every smooth member by smooth deformation
invariance of Gromov--Witten invariants, including the special quartic and
Gushel--Mukai presentations discussed in that appendix.
The general Fano cohomological shape and the positive $H^3$ dimensions in
genera two through five are those of the smooth classification
\cite[Introduction and Table~5]{KuznetsovProkhorovNodal}.
Lemmas~\ref{lem:cyclic-primary-persistence} and~\ref{lem:triple-primary-persistence}
continue the repeated factors in all even bulk directions;
Proposition~\ref{prop:rank2-rigidity} continues the rank-two residues.
By Lemma~\ref{lem:odd-selector}, no rank-one complement carries odd
cohomology. Hence all of $H^3$ belongs to the repeated factor, giving $O_3$.
Only its positivity is needed for irrationality.
\uses{prop:universal-rank-two-residue,lem:cyclic-primary-persistence,lem:triple-primary-persistence,prop:rank2-rigidity,lem:odd-selector}%
\imports{Fano:quantum-data,Fano:classification,Fano:reconstruction,Fano:genus-six-eight,Fano:gm-models,Przyjalkowski:counting-matrices}%
\evidence{fano-matrices}%
\end{proof}

\begin{proof}[Proof of Theorem~\ref{thm:fano-one-step-classification}]
The smooth classification has seventeen families
\cite[Introduction]{KuznetsovProkhorovNodal}.
The eight families outside $\mathcal F$ are rational: $\PP^3$, the smooth
quadric, index-two degrees four and five, and index-one genera seven,
nine, ten and twelve \cite[Section~1.1]{KuznetsovProkhorov}.
Their products with $\PP^1$ are rational.

For each of the remaining nine families, Proposition~\ref{prop:fano-endpoints}
gives $I_{\rm lat}>0$ or $O_3>0$.
Lemma~\ref{lem:special-projective-line} doubles the former.
For the latter, its tensor-product proof doubles the odd dimension at the
small point; a continued projector has constant rank on a formal germ,
so the same value holds in all even bulk directions.
Both counts vanish on surface centers and on $\PP^4$, and are birationally
invariant in dimension four. Thus these nine products are irrational.
A rational threefold has rational product with $\PP^1$, proving the reverse
implication as well. No deformation invariance of rationality is used.
\proves{thm:fano-one-step-classification}%
\uses{prop:fano-endpoints,lem:special-projective-line,lem:odd-selector,prop:numerical-blowup,prop:atomic-lowdim}%
\imports{Fano:classification,Fano:rational-controls,AKMW:weak-factorization}%
\end{proof}
